\section{JDBC}

\begin{question}
	以下关于Java中JDBC的说法错误的是
	\begin{options}
		\item 用于连接数据库
		\item 不同的数据库有不同的驱动
		\item 执行SQL语句通过Statement对象
		\item 不需要关闭连接资源
	\end{options}
	\begin{answer}
		D
	\end{answer}
	\begin{explanation}
		必须关闭 Connection、Statement、ResultSet 以释放资源。
	\end{explanation}
\end{question}

\begin{question}
	以下哪个方法用于执行查询语句？
	\begin{options}
		\item executeQuery()
		\item executeUpdate()
		\item execute()
		\item runQuery()
	\end{options}
	\begin{answer}
		A
	\end{answer}
	\begin{explanation}
		\texttt{executeQuery()} 用于 SELECT 查询，返回 ResultSet。
	\end{explanation}
\end{question}

\begin{question}
	以下关于Java中PreparedStatement的说法正确的是
	\begin{options}
		\item 性能比Statement差
		\item 不能防止SQL注入
		\item 可以设置参数
		\item 不支持批量操作
	\end{options}
	\begin{answer}
		C
	\end{answer}
	\begin{explanation}
		PreparedStatement 可预编译，支持参数化查询，防止 SQL 注入，性能更好。
	\end{explanation}
\end{question}

\begin{question}
	以下关于Java中事务的说法错误的是
	\begin{options}
		\item 保证数据的一致性
		\item 可以通过commit()提交
		\item 可以通过rollback()回滚
		\item 事务中的操作一定能成功
	\end{options}
	\begin{answer}
		D
	\end{answer}
	\begin{explanation}
		事务中的操作可能失败，此时应回滚。
	\end{explanation}
\end{question}
